\section{Methods}\label{sec:methods}

\subsection{Data-fitted neuron models}
We used the differentiable generalized firing-rate (GFR) models of \citet{yeung2025gfr}.
They were fitted to whole-cell current-clamp recordings in the Allen Cell Types Database
\citep{gouwens2019}, and we took the code and data from upstream commit \texttt{835a535}. Each
cell $i$ has $K=8$ filters with fixed time constants
$\tau_k\in\{{\approx}0, 20, 50, 100, 200, 500, 1000, 2000\}$\,ms. At the fitted bin size
$\Delta t=20$\,ms the decay coefficients are $d_k = 1-e^{-\Delta t/\tau_k}$, shared by all cells.
With input current $I_{i,t}$ (pA) and rate $r_{i,t}$ (Hz),
\begin{align}
h^{(k)}_{i,t} &= (1-d_k)\,h^{(k)}_{i,t-1} + a_{ik}\,I_{i,t} + \beta_{ik}\,r_{i,t-1},\label{eq:gfr-h}\\
r_{i,t} &= \gamma_i\,\mathrm{ReLU}\!\Big[\tanh\Big(\textstyle\sum_{d}c_{id}^2\,x_{i,t}^{\,d}\Big)\Big],
\qquad x_{i,t}=\big(\bar h_{i,t}-b_i\big)/\sigma_i,\quad \bar h_{i,t}=\tfrac1K\textstyle\sum_k h^{(k)}_{i,t}.\label{eq:gfr-r}
\end{align}
The stored $a_{ik}$ and $\beta_{ik}$ already include the discretization, so no extra $\Delta t$
factor is applied. Upstream rates in spikes/ms are converted to Hz once. Fitted parameters are
never changed by training. Our implementation reproduces the upstream equation for all 1,124
fitted cells (unit tests).

The coefficients $\beta_{ik}$ form the cell's spike-history filter. Its steady-state gain
$\bar\beta_i=\tfrac1K\sum_k\beta_{ik}/d_k$ (pA\,Hz$^{-1}$) is negative for almost every cell,
which is spike-frequency adaptation \citep{benda2003,weber2017glm}: median $-10.5$ for pyramidal
(Pyr) cells, $-6.6$ for SST, $-4.1$ for VIP and $-0.8$ for PV cells in the development cohort.

\paragraph{Cells and cohorts.}
Following the upstream release criterion, we kept mouse cells with validation explained
variance $\ge0.5$. We did not select on the held-out test sweeps. Operational labels came from
Cre line and dendrite type: spiny cells from excitatory lines are Pyr, and Pvalb-, Sst- and
Vip-Cre cells are PV, SST and VIP. This gave 597 cells (339 Pyr, 107 PV, 86 SST, 65 VIP) from 410
donors, mostly in primary visual cortex. Whole donors were assigned to a development cohort (298
cells) or a confirmation cohort (299 cells). The assignment was a seeded, type-stratified greedy
procedure. No donor contributes to both cohorts. All tuning and piloting used only the development
cohort.

\subsection{Networks}
A network has $N=100$ cells, 80 Pyr, 10 PV, 5 SST and 5 VIP, sampled without replacement for each
network seed. The input current to cell $i$ is
\begin{equation}
I_{i,t}=I^0_i + S_i\Big(\textstyle\sum_c W^{\rm in}_{ic}u_{c,t}+\sum_j s_j |W_{ij}| M_{ij}\, r_{j,t-1}/R_0\Big),
\end{equation}
with $R_0=100$\,Hz for every cell. The sign $s_j$ is $+1$ for Pyr and $-1$ for interneurons
(Dale's law). The mask $M$ is drawn from a type-level connection-probability matrix inspired by
cortical motifs \citep{pfeffer2013,pi2013}. These are modelling assumptions, not measured values
(Table~\ref{tab:priors}). Magnitudes $|W_{ij}|$ are log-normal with mean 1 ($\sigma_{\log}=0.35$),
scaled by $0.6/\langle\text{in-degree}\rangle$, and inhibitory sources are multiplied by 1.5.
Input weights are $W^{\rm in}_{ic}\sim\mathcal N(0,1/5)$. The readout
$z_t=\sum_{i\in\rm Pyr}w_i r_{i,t}/R_0+b$ uses Pyr cells only.

\paragraph{Operating-point calibration.}
For each cell, $I^0_i$ is chosen on a 33-point grid in $[0,\sigma_i]$ so that its isolated
steady-state rate is closest to 5\% of its own fitted ceiling $\gamma_i$. The rate is measured in
the last third of a 3-s constant-current probe. The input scale is $S_i=0.25\,\sigma_i$. The
procedure uses no task information. Because the target is relative to each cell's ceiling, the
resulting baseline rates differ by type (PV higher), by construction.

\paragraph{Identity and wiring controls.}
\emph{Within-type shuffle}: independent permutations, within each type, of the input filters
$a_{i\cdot}$, the history filters $\beta_{i\cdot}$ and the activation tuple
$(\sigma,\gamma,c,b)$. This keeps marginal distributions and breaks the pairing of a cell's own
parameters. \emph{Type-homogeneous}: every parameter is replaced by its type mean.
\emph{Cross-type permutation}: whole parameter tuples are permuted across types. \emph{Fast
decay}: $d_k\leftarrow\max(d_k,1-e^{-1})$, so $\tau_k\le20$\,ms, with $a$ and $\beta$ rescaled to
keep $a/d$ and $\beta/d$, and hence the static response. Unless stated otherwise, each manipulated
bank is recalibrated with the same label-free procedure (``matched''). The ``reference''
calibration keeps the intact bank's $I^0,S$, as in earlier versions of this project; it is shown as
a control. \emph{Degree-preserving rewiring} applies $10|E|$ directed double-edge swaps. It
preserves every cell's in- and out-degree, its Dale sign and its multiset of outgoing magnitudes.

\subsection{Tasks}
Inputs have five channels: two sensory streams (A, B), two cue channels and a go signal. Time is
binned at 20\,ms: fixation 100\,ms, stimulus 200\,ms, a silent delay of 100, 200 or 300\,ms in
training, and a 60-ms response window marked by the go channel. Sensory evidence has signed
coherence $\pm\{0.25,0.5,1\}$ plus Gaussian noise ($\sigma=0.25$ per bin). Every batch is exactly
balanced over cue, relevant sign and irrelevant sign.
\emph{T1-early} \citep[context-dependent decision;][]{mante2013}: a context cue during fixation and
stimulus says which stream is relevant, and the target is the sign of the relevant stream.
\emph{T1-late}: the same, but the context cue (60\,ms) arrives only after the delay, so both
streams must be held. \emph{T2-dms} (delayed match-to-sample, \citealp{miller1996,masse2019}):
stream A is the sample and stream B a simultaneous distractor. A $\pm1$ probe arrives after the
delay, and the target is $\mathrm{sign(sample)}\times\mathrm{probe}$, so no input before the probe
carries information about the choice. Conflict trials are those whose two streams have opposite
signs. Test splits are 512 trials each: in-distribution (ID); unseen delays of 400 and 600\,ms
(OOD); higher sensory noise ($\sigma=0.5$); a 40-ms zero-mean distractor ($\sigma=0.75$) in
mid-delay; and all three combined.

\subsection{Loss and learning rules}
The loss is the binary cross-entropy of $z_t$ at every response bin, averaged within a trial and
then over the batch. The choice is the sign of the mean response-window logit. Only $|W|$ (masked,
non-negative, at most 2), $w$ and $b$ are plastic.

\paragraph{BPTT.} The exact gradient by back-propagation through time \citep{werbos1990bptt}
through all filters, the spike-history loop and the recurrent connections. It was checked against
central finite differences.

\paragraph{Local rule.} The local rule is e-prop \citep{bellec2020eprop}, which is related to
RFLO \citep{murray2019rflo}, adapted to GFR cells. It approximates $\partial\mathcal L/\partial
W_{ij}$ by
\begin{equation}
\hat g_{ij}=\sum_t L_{i,t}\,\frac{g'_{i,t}}{R_0}\,\bar e_{ij,t},\qquad
L_{i,t}=B_i\,\frac{\partial \mathcal L}{\partial z_t},
\end{equation}
where $B$ is a fixed random feedback vector \citep{lillicrap2016fa} and $g'$ is the activation
slope. The $K$-dimensional eligibility trace $e_{ij,t}=\partial h_{i,t}/\partial W_{ij}$ follows
only cell $i$'s own dynamics, with presynaptic activity held fixed:
\begin{equation}
e^{(k)}_{ij,t}=(1-d_k)\,e^{(k)}_{ij,t-1}+a_{ik}S_i s_j\,\frac{r_{j,t-1}}{R_0}
+\underbrace{\beta_{ik}\,g'_{i,t-1}\,\bar e_{ij,t-1}}_{\text{spike-history term}}.\label{eq:elig}
\end{equation}
Every quantity is available at the synapse or the postsynaptic cell, plus one broadcast error
signal \citep{gerstner2018}. \emph{No-$\beta$} removes only the spike-history term in
Eq.~\ref{eq:elig}; the network dynamics are unchanged. \emph{Symmetric feedback} uses $B=w$ (the
readout weights), which isolates the error introduced by ignoring paths through other cells. The
quantity in Eq.~\ref{eq:elig} can be computed forward in time, as a synapse would, or by an
equivalent reverse sweep over each cell's own dynamics, which is the network gradient with all
paths through other cells removed. The two agree to $<10^{-10}$ relative error (unit tests). We
used the reverse sweep because it is $\sim$17$\times$ faster.

\paragraph{Frozen reservoir.} Recurrent weights are fixed, and only $w,b$ are trained, with their
exact gradient.

\paragraph{Optimizers.} \emph{Adam} \citep{kingma2015adam} with global gradient-norm clipping at
1.0 is the primary optimizer. \emph{RMS} is per-parameter RMS normalization without momentum and
without clipping \citep{tieleman2012rmsprop}, so every update depends only on that synapse's own
history. It is an optimizer, not a model of homeostatic synaptic scaling \citep{turrigiano2008}.
After each update, magnitudes are projected onto $[0,2]$.

\paragraph{Training and model selection.} Each run trains for 3000 updates of 32 trials.
Validation uses 256 ID trials every 50 updates, and the checkpoint with the lowest validation BCE
(including the untrained network) is kept. Learning rates were chosen on the development cohort
only. The grid was $\eta_{\rm rec}\in\{3\cdot10^{-4},10^{-3},3\cdot10^{-3},10^{-2}\}\times
\eta_{\rm out}\in\{0.03,0.1,0.3,1\}$, three seeds per candidate, and the criterion was mean best
validation BCE. This was done separately for each learner, task and optimizer
(Table~\ref{tab:lr}). A second stage tried $\eta_{\rm out}\in\{3,10\}$ at each learner's
selected $\eta_{\rm rec}$, because with Adam the frozen reservoir and the no-$\beta$ rule were
still improving at the top of the grid. The final choice minimized the same criterion over both
stages; the reservoir's optimum ($\eta_{\rm out}=3$) is then inside the range tried. An earlier
search with $\eta_{\rm out}\le0.1$ had put every learner at its upper edge and was discarded. No
test split had been evaluated at any of these points. Controls use the rates of their learner. A network seed fixes the sampled
cells, graph, initial weights, feedback vector and every trial, so all conditions are paired.

\subsection{Analyses}
\paragraph{Population codes.} For each phase (stimulus, delay, response), we fitted PCA, ridge
probes ($\alpha=1$) for choice, cue and each stream's sign, and a class-mean choice axis on 128
\emph{analysis-fit} trials. All were evaluated on 128 independent \emph{analysis-evaluation}
trials, to avoid double dipping \citep{kriegeskorte2009}.

\paragraph{Gradient alignment.} At each saved checkpoint, on 32 independent diagnostic trials, we
computed the cosine between the exact recurrent gradient and each local estimate: with or without
$\beta$, and with random or symmetric feedback. Per-neuron alignment uses the row of the gradient
belonging to postsynaptic cell $i$. Each cell's \emph{adaptation loop gain} is
$|\langle g'_i\rangle\,\bar\beta_i|$, the steady-state history gain times the mean activation
slope on those trials.

\subsection{Statistics and protocol}
The unit of inference is the network. We avoid pooling cells or trials as if they were
independent \citep{aarts2014}. Comparisons are paired by seed and use exact two-sided sign-flip
tests over all $2^{16}$ sign assignments, with percentile bootstrap CIs over networks. Families of
tests and the Holm correction are as preregistered (Section~\ref{sec:protocol}). The 2$\times$2$\times$2
factorial effects (identity, wiring, learner) are main effects and interactions computed within
each network. An identity, wiring or heterogeneity effect is called \emph{practically negligible}
when its 90\% CI lies within $\pm5$\pp. The protocol, the analysis code and the source hash were
locked before the confirmation cohort was used (Supplementary Section~\ref{sec:protocol}). Results
outside it are labelled exploratory.
